\begin{abstract}
This paper tests seven candidate interventions for peak hour congestion in Kathmandu on a
mesoscopic simulation whose demand is built from the 2011 JICA vehicle and person trip origin
destination survey, on a network of arterial and tertiary roads with side streets that give way
at junctions. Road speeds are calibrated to the JICA travel speed survey only; intersection counts
are held out for validation, where the model carries \ValShare{} of counted peak volume with a
junction correlation of \ValCorr. Every comparison uses 15 paired random seeds and mean time in
system over all vehicles, so trips that never finish cannot flatter a result. Staggered departure
times reduce time in system by \HSixTisMag\%, adaptive signal control by \HOneTisMag\%, and
adaptive signals with congestion pricing by \HThreeTisMag\%, all well beyond their 95\% intervals.
Pricing alone gives \HTwoTisMag\%, and the combined effect is the sum of its parts, not more.
Decentralizing trip attraction helps only as much as cutting the same volume uniformly; tidal
lanes and a parking enforcement proxy show no reliable effect. Under demand grown to 2030, only
\ProjThirtyComp{} of peak trips finish. Compared with an earlier version of this model on synthetic
demand and main roads only, the ranking of interventions holds but pricing loses most of its
effect, a reminder that intervention estimates depend on how much of the street network the model
represents.

\medskip\noindent\textbf{Keywords:} traffic simulation, signal control, congestion pricing, travel
demand management, Kathmandu, mesoscopic modeling
\end{abstract}
